% !TEX root = ../main.tex
% Figure: Dynamics and modulation response of PCSEL
% Chapter: ch21b - Dynamics, Modulation, and Noise Characteristics
% Description: Small-signal modulation bandwidth and resonance peak

\begin{figure}[htbp]
  \centering
  \begin{tikzpicture}[scale=0.88]
    % AM frequency response plot
    \begin{scope}
      % Axes
      \draw[->, thick] (0,0) -- (7,0) node[right] {频率$f$（千兆赫）};
      \draw[->, thick] (0,-4) -- (0,3) node[above] {归一化响应（分贝）};
      
      % Grid lines
      \foreach \y in {-3,-2,...,2} {
        \draw[gray!30] (0,\y) -- (6.5,\y);
      }
      
      % Three curves for different bias currents
      \draw[very thick, blue, domain=0:6.5, samples=100] 
        plot(\x,{max(-4, 20*log10(1/sqrt((1-(\x/2)^2)^2+(0.5*\x/2)^2)))-3});
      \node[blue, above] at (3,1) {\scriptsize $I/I_{th}=1.5$};
      
      \draw[very thick, red, domain=0:6.5, samples=100] 
        plot(\x,{max(-4, 20*log10(1/sqrt((1-(\x/3)^2)^2+(0.4*\x/3)^2)))-1});
      \node[red, above] at (4,1.5) {\scriptsize $I/I_{th}=2.5$};
      
      \draw[very thick, green!60!black, domain=0:6.5, samples=100] 
        plot(\x,{max(-4, 20*log10(1/sqrt((1-(\x/4)^2)^2+(0.3*\x/4)^2)))+1});
      \node[green!60!black, above] at (5,2) {\scriptsize $I/I_{th}=4.0$};
      
      % Resonance peaks marking
      \fill[blue] (1.9,-0.5) circle (2pt);
      \fill[red] (2.9,0.3) circle (2pt);
      \fill[green!60!black] (3.9,1.2) circle (2pt);
      
      % -3dB bandwidth indication
      \draw[dashed, thick] (0,-3) -- (5.5,-3);
      \draw[<->, thick, magenta] (4.2,-3.3) -- node[below] {\scriptsize $f_{-3\,\text{分贝}}$} (5.5,-3.3);
      
      % Relaxation oscillation frequency
      \draw[dotted, thick] (2.9,0) -- (2.9,3);
      \node[below, red] at (2.9,-0.3) {$f_R$};
      
      \node[above, font=\bfseries] at (3.5,3.5) {(a) 小信号调制响应};
    \end{scope}
    
    % Eye diagram inset
    \begin{scope}[shift={(8,0)}, scale=0.8]
      % Axes
      \draw[->, thick] (0,0) -- (5,0) node[right] {时间（单位间隔）};
      \draw[->, thick] (0,-2) -- (0,2.5) node[above] {光功率};
      
      % Ideal eye pattern
      \draw[very thick, blue, opacity=0.6] 
        (0,0) .. controls (1,0.5) and (1.5,1.5) .. (2.5,2)
        (0,2) .. controls (1,1.5) and (1.5,0.5) .. (2.5,0)
        (2.5,0) .. controls (3.5,-0.5) and (4,0.5) .. (5,2)
        (2.5,2) .. controls (3.5,1.5) and (4,0.5) .. (5,0);
      
      % Real degraded pattern (overlaid)
      \draw[thick, red, opacity=0.4] 
        (0,0.2) .. controls (1,0.6) and (1.5,1.3) .. (2.5,1.8)
        (0,1.8) .. controls (1,1.4) and (1.5,0.7) .. (2.5,0.2);
      
      % Eye opening measurement
      \draw[<->, thick, green!60!black] (2.3,-0.3) -- (2.3,0.3);
      \node[right, green!60!black] at (2.4,0) {\scriptsize 眼高};
      
      \draw[<->, thick, orange] (1.8,1) -- (3.2,1);
      \node[above, orange] at (2.5,1.1) {\scriptsize 眼宽};
      
      \node[above, font=\bfseries] at (2.5,2.8) {(b) 眼图};
      
      % Quality metrics
      \node[align=left, font=\footnotesize, anchor=north west] at (-0.5,-2.5) {
        \textbf{信号质量指标}:\\
        $\bullet$ 消光比: $P_1/P_0 > 8\,\text{分贝}$\\
        $\bullet$ 抖动: $<0.1\,\text{单位间隔}$\\
        $\bullet$ Q 因子:$>6$ （误码率$<10^{-9}$）
      };
    \end{scope}
  \end{tikzpicture}
  \caption{光子晶体面发射激光器动态调制示意：左图为小信号频率响应随偏置电流升高而展宽，右图为高速数字调制下的眼图与主要质量指标。}
  \label{fig:ch21b_modulation_response}
\end{figure}
